\chapter{Chains and local Hamiltonians}
\label{ch:one_dimensional_preliminaries}

This chapter collects the definitions and elementary facts about chains
used by the one-dimensional chapters that follow: the chain of $L$ sites
with local dimension $q$, its intervals, finite-range Hamiltonians with one
bounded Hermitian term per support, counts of the supports that cross a
cut, and Lieb--Robinson bounds on a contiguous chain. The area law itself
is proved in Chapter~\ref{ch:one_dimensional_chain}.

\section{Chains, supports and Hamiltonians}

\begin{definition}[Chains, intervals and finite-range Hamiltonians]
    \label{def:chain_model}
    \lean{
        TNLean.MPS.AreaLaw.interval,
        TNLean.MPS.AreaLaw.mem_interval,
        TNLean.MPS.AreaLaw.interval_mono,
        TNLean.MPS.AreaLaw.interval_eq_empty,
        TNLean.MPS.AreaLaw.interval_zero_eq_univ,
        TNLean.MPS.AreaLaw.card_interval_le,
        TNLean.MPS.AreaLaw.IsChainSupport,
        TNLean.MPS.AreaLaw.ChainSupport,
        TNLean.MPS.AreaLaw.IsChainSupport.card_le,
        TNLean.MPS.AreaLaw.supportDim_le,
        TNLean.MPS.AreaLaw.one_le_supportDim,
        TNLean.MPS.AreaLaw.ChainHamiltonian,
        TNLean.MPS.AreaLaw.ChainHamiltonian.operator,
        TNLean.MPS.AreaLaw.ChainHamiltonian.operator_isHermitian,
        TNLean.MPS.AreaLaw.ChainHamiltonian.zero,
        TNLean.MPS.AreaLaw.cutLogBudget,
        TNLean.MPS.AreaLaw.one_le_cutLogBudget
    }
    \leanok
    The \emph{chain of $L$ sites} has sites $0,\ldots,L-1$, each carrying
    $\mathbb C^q$. For natural numbers $a\le b$ the \emph{interval} $[a,b)$ is
    the set of sites with position at least $a$ and less than $b$; it has at
    most $b-a$ sites, and its dilation $[a-d,b+d)$ is clipped at the ends of the
    chain. A \emph{support of range $R$} is a nonempty set of sites of diameter
    at most $R$; it has at most $R+1$ sites and Hilbert-space dimension at most
    $q^{R+1}$. A \emph{finite-range Hamiltonian} assigns to every support $X$
    of range $R$ a Hermitian term $h_X$ acting as the identity off $X$ with
    operator norm at most $J$, zero terms allowed, and
    $H=\sum_X h_X$. A \emph{gapped ground vector} with gap $\Delta$ is a unit
    vector $\Omega$ with $H\Omega=E_0\Omega$ and
    $H-E_0\ge\Delta(1-|\Omega\rangle\langle\Omega|)$. For a cut $A$ the
    \emph{cut budget} is
    $\mathcal B_A=1+\sum_X\log^2(e\,q^{|X|})$ over the supports $X$ meeting
    both $A$ and its complement; it is at least one.
\end{definition}

\begin{lemma}[A bond sum separated by two windows]
    \label{thm:chain_two_window_sum}
    \lean{Finset.sum_range_eq_five_parts_of_two_windows,
        Fin.sum_univ_eq_five_parts_of_two_windows}
    \leanok
    Let $M$ be an abelian monoid written additively, and let
    $a,b,\ell,N$ be natural numbers with $a+2\ell\le b$ and
    $b+2\ell\le N$. For a sequence $f:\mathbb N\to M$,
    \begin{align}
        \sum_{t<N}f(t)
         & =\sum_{t<a}f(t)+\sum_{r<2\ell}f(a+r)\notag   \\
         & \quad+\sum_{r<b-(a+2\ell)}f(a+2\ell+r)\notag \\
         & \quad+\sum_{r<2\ell}f(b+r)
        +\sum_{r<N-(b+2\ell)}f(b+2\ell+r).
        \label{eq:chain_two_window_sum}
    \end{align}
    The same identity holds when $f$ is given only on
    $\{0,\ldots,N-1\}$: every shifted index in~(\ref{eq:chain_two_window_sum})
    belongs to this set. Empty windows and empty complementary intervals
    are permitted. Applied to a bond array, this is the additive grouping
    used in~\cite[Section~4.1]{Arad2016RigorousRG}.
\end{lemma}
\begin{proof}\leanok
    Split the sum successively at $b+2\ell$, $b$, $a+2\ell$ and $a$.
    The inequalities on the endpoints give the complementary lengths
    in~(\ref{eq:chain_two_window_sum}). For a finite array, extend $f$ by zero
    outside its domain and reindex each interval by its finite length.
    Each shifted index is less than $N$, so the extension agrees with $f$
    in every summand.
\end{proof}

\begin{definition}[Two-window grouping of a bond family]
    \label{def:chain_two_window_family}
    \lean{Fin.twoWindowFamily}
    \leanok
    Let $A$ be an additive commutative monoid, and let
    $h_0,\ldots,h_{N-1}\in A$. For the two-window decomposition of
    \cite{Arad2016RigorousRG}, take natural numbers $a,b,\ell$ with
    \begin{align}
        a+2\ell & \le b,\notag \\
        b+2\ell & \le N.
        \label{eq:chain_two_window_admissibility}
    \end{align}
    Put $W=\{r\in\mathbb N:r<2\ell\}$ and form the disjoint union
    \begin{align}
        I & =(W\sqcup W)\sqcup\{L,M,R\}.
        \label{eq:chain_two_window_family_alphabet}
    \end{align}
    Write $\iota_L,\iota_R$ for the two window injections and
    $\kappa$ for the bulk injection. Thus $|I|=4\ell+3$, and the two
    window injections have disjoint images. Define $F:I\to A$ by
    \begin{align}
        F_{\iota_L(r)} & =h_{a+r},\notag                  \\
        F_{\iota_R(r)} & =h_{b+r},\notag                  \\
        F_{\kappa(L)}  & =\sum_{0\le t<a}h_t,\notag       \\
        F_{\kappa(M)}  & =\sum_{a+2\ell\le t<b}h_t,\notag \\
        F_{\kappa(R)}  & =\sum_{b+2\ell\le t<N}h_t.
        \label{eq:chain_two_window_family_values}
    \end{align}
    The three bulk labels remain distinct when their sums vanish.
\end{definition}

\begin{lemma}[Sum of the grouped bond family]
    \label{thm:chain_two_window_family_sum}
    \lean{Fin.sum_twoWindowFamily}
    \leanok
    \uses{def:chain_two_window_family}
    The family in Definition~\ref{def:chain_two_window_family} satisfies
    \begin{align}
        \sum_{\alpha\in I}F_\alpha & =\sum_{t=0}^{N-1}h_t.
        \label{eq:chain_two_window_family_sum}
    \end{align}
\end{lemma}
\begin{proof}\leanok
    \uses{def:chain_two_window_family,thm:chain_two_window_sum}
    Put $S_L=\sum_{r\in W}h_{a+r}$, $S_R=\sum_{r\in W}h_{b+r}$ and
    $B_\nu=F_{\kappa(\nu)}$ for $\nu\in\{L,M,R\}$.
    Expanding~(\ref{eq:chain_two_window_family_values}) gives
    \begin{align}
        \sum_{\alpha\in I}F_\alpha
         & =(S_L+S_R)+(B_L+B_M+B_R)\notag \\
         & =B_L+S_L+B_M+S_R+B_R.
        \label{eq:chain_two_window_regrouping}
    \end{align}
    Lemma~\ref{thm:chain_two_window_sum} identifies the expression
    in~(\ref{eq:chain_two_window_regrouping}) with the original sum,
    proving~(\ref{eq:chain_two_window_family_sum}). Empty intervals
    cause no exception.
\end{proof}

\begin{lemma}[Additive maps and the grouped bond family]
    \label{thm:chain_two_window_family_map}
    \lean{Fin.twoWindowFamily_map}
    \leanok
    \uses{def:chain_two_window_family}
    Let $A,B$ be additive commutative monoids, let $\varphi:A\to B$ be a
    homomorphism of additive monoids, and let $F:I\to A$ be the family in
    Definition~\ref{def:chain_two_window_family}. Form $F':I\to B$ from the
    array $\varphi(h_0),\ldots,\varphi(h_{N-1})$ using the same windows.
    Then $F'_\alpha=\varphi(F_\alpha)$ for every $\alpha\in I$.
\end{lemma}
\begin{proof}\leanok
    \uses{def:chain_two_window_family}
    For a window label the assertion follows from the definition.
    For each of the three bulk labels, apply preservation of finite sums
    by $\varphi$ to the corresponding complementary interval.
\end{proof}

\begin{lemma}[Three-region tensor forms after grouping two windows]
    \label{thm:chain_two_window_tensor_forms}
    \lean{Matrix.twoWindowFamily_eq_three_region_forms}
    \leanok
    \uses{def:chain_two_window_family}
    Let $X,Y,Z$ be index sets, and let $H_r$, $H^L_r$, $H^M_r$, $H^R_r$
    be complex matrices indexed respectively by $(X\times Y)\times Z$,
    $X$, $Y$, $Z$, for $0\le r<N$.
    Use the windows and alphabet of Definition~\ref{def:chain_two_window_family},
    and choose $p,s\in W$. For matrices $A,B,C$ on the three factors, put
    \begin{align}
        \mathcal E_L(A) & =(A\otimes I_Y)\otimes I_Z,\notag \\
        \mathcal E_M(B) & =(I_X\otimes B)\otimes I_Z,\notag \\
        \mathcal E_R(C) & =I_{X\times Y}\otimes C.
        \label{eq:chain_two_window_embeddings}
    \end{align}
    Suppose that, whenever $r\ne a+p$ and $r\ne b+s$,
    \begin{align}
        H_r & =\mathcal E_L(H^L_r)+\mathcal E_M(H^M_r)+\mathcal E_R(H^R_r).
        \label{eq:chain_two_window_family_ordinary_forms}
    \end{align}
    Group each of the four arrays using the same windows, obtaining
    $F,F^L,F^M,F^R$. Then, for every $\alpha\in I$ other than
    $\iota_L(p)$ and $\iota_R(s)$,
    \begin{align}
        F_\alpha & =\mathcal E_L(F^L_\alpha)+\mathcal E_M(F^M_\alpha)
        +\mathcal E_R(F^R_\alpha).
        \label{eq:chain_two_window_grouped_forms}
    \end{align}
    This is the additive step in the grouping used in~\cite{Arad2016RigorousRG}; the original tensor identities are hypotheses
    of the assertion. Empty index sets are allowed.
\end{lemma}
\begin{proof}\leanok
    \uses{def:chain_two_window_family}
    Each map in~(\ref{eq:chain_two_window_embeddings}) is linear. Hence, for
    any finite set $J$ of original indices avoiding $a+p$ and $b+s$,
    \begin{align}
        \sum_{r\in J}H_r
         & =\mathcal E_L\left(\sum_{r\in J}H^L_r\right)
        +\mathcal E_M\left(\sum_{r\in J}H^M_r\right)
        +\mathcal E_R\left(\sum_{r\in J}H^R_r\right).
        \label{eq:chain_two_window_sum_forms}
    \end{align}
    For a window label different from the two selected labels,
    use~(\ref{eq:chain_two_window_family_ordinary_forms}). Each complementary
    interval avoids both selected indices by
    (\ref{eq:chain_two_window_admissibility}). Applying
    (\ref{eq:chain_two_window_sum_forms}) to each of these three intervals
    proves~(\ref{eq:chain_two_window_grouped_forms}) for the bulk labels.
\end{proof}

\begin{lemma}[The number of range-one supports]
    \label{thm:chain_range_one_support_count}
    \lean{
        TNLean.MPS.AreaLaw.card_chainSupport_range_one_le,
        TNLean.MPS.AreaLaw.ChainHamiltonian.norm_operator_le_of_range_one
    }
    \leanok
    \uses{def:chain_model}
    A chain of $L$ sites has at most $2L$ supports of range one.
    Consequently, a range-one Hamiltonian whose terms have norm at most
    $J$ satisfies $\|H\|\le2LJ$.
\end{lemma}
\begin{proof}\leanok
    A support of range one is a singleton or a pair of adjacent sites.
    It is therefore determined by its least site and its cardinality,
    which is either one or two. This gives the bound $2L$.
    If $L>0$, the norm bound on a singleton term implies $J\ge0$, and the
    triangle inequality gives $\|H\|\le\sum_X\|h_X\|\le2LJ$.
    If $L=0$, there are no interaction terms and both sides vanish.
\end{proof}

The following upper bound on the gap is an auxiliary estimate for the constant
comparison in~\cite[lines 291--293 and 509--515]{Landau2015Polynomial}.

\begin{lemma}[An upper bound on the gap of a range-one chain]
    \label{thm:chain_range_one_gap_upper}
    \lean{TNLean.MPS.AreaLaw.ChainHamiltonian.gap_le_four_mul_of_range_one}
    \leanok
    \uses{def:chain_model}
    Let $L\ge1$, $q\ge2$ and $J\in\mathbb R$, and let $H=\sum_Xh_X$ be a
    range-one Hamiltonian on the chain with $\lVert h_X\rVert\le J$.
    Let $\Omega$ be a gapped ground vector of $H$, with gap
    $\Delta\in\mathbb R$. Then $\Delta\le4J$.
\end{lemma}

\begin{proof}
    \leanok
    \uses{def:chain_model}
    Diagonalize the marginal $\rho_0$ at site $0$, and let $U_0$ cyclically
    permute an orthonormal eigenbasis. Since $q\ge2$, this permutation has
    no fixed point, so $\operatorname{tr}(\rho_0U_0)=0$. Put $K=U_0\otimes I$
    and $\varphi=K\Omega$. Since $K$ acts only on site $0$ and is unitary,
    \begin{align}
        \langle\Omega,\varphi\rangle & =\langle\Omega,K\Omega\rangle
        =\operatorname{tr}(\rho_0U_0)=0,\notag                       \\
        \lVert\varphi\rVert          & =\lVert K\Omega\rVert
        =\lVert\Omega\rVert=1.\notag
    \end{align}
    If $0\notin X$, then $h_X$ acts as the identity on site $0$, so
    $K^\dagger h_XK=h_X$ and
    \begin{align}
        \langle\varphi,h_X\varphi\rangle
        =\langle K\Omega,h_XK\Omega\rangle
        =\langle\Omega,K^\dagger h_XK\Omega\rangle
        =\langle\Omega,h_X\Omega\rangle.\notag
    \end{align}
    At most two range-one supports contain $0$: the singleton $\{0\}$ and,
    when it exists, the edge $\{0,1\}$. Each expectation changes by at most
    $2J$. The gap inequality gives
    \begin{align}
        \Delta\le\langle\varphi,H\varphi\rangle-E_0
        =\sum_{X\ni0}\bigl(\langle\varphi,h_X\varphi\rangle
        -\langle\Omega,h_X\Omega\rangle\bigr)\le4J.\notag
    \end{align}
\end{proof}

\begin{lemma}[A range-one support crossing a cut]
    \label{thm:chain_range_one_crossing_support}
    \lean{TNLean.MPS.AreaLaw.IsChainSupport.eq_interval_of_splits_of_range_one}
    \leanok
    \uses{def:chain_model}
    If a support $X$ of range one contains sites on both sides of the cut
    after the first $k$ sites, then $X=\lbrack k-1,k+1)$.
\end{lemma}
\begin{proof}\leanok
    \uses{thm:chain_full_support_window}
    The support contains distinct sites $x<k\le y$, so it has at least two
    sites. Its diameter is at most one, so it has exactly two sites. By
    Lemma~\ref{thm:chain_full_support_window} with $R=1$, it is an interval
    $\lbrack a,a+2)$ of two adjacent sites, and $a\le x<k\le y\le a+1$ forces
    $a=k-1$.
\end{proof}

\begin{lemma}[The Hamiltonian at one cut]
    \label{thm:chain_range_one_cut_decomposition}
    \lean{TNLean.MPS.AreaLaw.ChainHamiltonian.exists_cut_decomposition_of_range_one}
    \leanok
    \uses{def:chain_model}
    Let $H=\sum_Xh_X$ be a range-one Hamiltonian whose terms are positive
    semidefinite and satisfy $\|h_X\|\le1$. For every natural number $k$,
    let $A=\lbrack0,k)$ be the interval clipped to the chain. There is a
    decomposition
    \begin{align}
        H & =H_L+H_R+K,\notag \\
        0 & \le K\le I.\notag
    \end{align}
    where $H_L$ is supported on $A$ and $H_R$ on its complement.
    This gives the decomposition used for nearest-neighbour interactions
    in~\cite[lines 283 and 409--410]{Landau2015Polynomial}; positive
    one-site terms are also allowed here.
\end{lemma}
\begin{proof}\leanok
    \uses{thm:chain_range_one_crossing_support}
    Let $C$ be the set of interaction supports meeting both $A$ and $A^{\mathsf c}$.
    Partitioning the nonempty interaction supports gives
    \begin{align}
        H & =\underbrace{\sum_{X\subseteq A}h_X}_{H_L}
        +\underbrace{\sum_{X\subseteq A^{\mathsf c}}h_X}_{H_R}
        +\underbrace{\sum_{X\in C}h_X}_{K}.\notag
    \end{align}
    The first two sums have the required supports. By
    Lemma~\ref{thm:chain_range_one_crossing_support}, $|C|\le1$.
    Thus $K$ is either zero or one positive semidefinite term of norm at
    most one, and $0\le K\le I$.
\end{proof}

\begin{lemma}[Decomposition for boundary gluing]
    \label{thm:chain_gluing_hamiltonian_decomposition}
    \lean{TNLean.MPS.AreaLaw.ChainHamiltonian.exists_gluing_decomposition_of_range_one}
    \leanok
    \uses{def:chain_model}
    Let $H=\sum_Xh_X$ be a range-one Hamiltonian on $L$ sites with
    $\lVert h_X\rVert\le1$. For every natural number $k$, put
    $A=\lbrack0,k)$ and $D=\lbrack0,\max\{k-1,0\})^{\mathsf c}$, with intervals clipped
    to the chain. There exist matrices $H_L$ and $B$ such that
    \begin{align}
        H              & =H_L+B,\notag \\
        \lVert B\rVert & \le2L.\notag
    \end{align}
    where $H_L$ is supported on $A$ and $B$ on $D$.
    This supplies the decomposition $H=H_L+(H_R+H_i)$ used
    in~\cite[lines 409--410]{Landau2015Polynomial} for the range-one model
    of Definition~\ref{def:chain_model}, which also allows Hermitian one-site terms.
\end{lemma}
\begin{proof}\leanok
    \uses{thm:chain_range_one_support_count}
    Put $H_L=\sum_{X\subseteq A}h_X$ and
    $B=\sum_{X\nsubseteq A}h_X$. Then $H=H_L+B$, and $H_L$ is
    supported on $A$. Every support in the second sum contains a site
    $y\ge k$. For each of its sites $x$, the range-one condition gives
    $y\le x+1$, so $k\le x+1$ and $x\in D$. Thus $B$ is supported on $D$.
    By Lemma~\ref{thm:chain_range_one_support_count}, there are at most
    $2L$ supports. The triangle inequality and $\lVert h_X\rVert\le1$ give
    \begin{align}
        \lVert B\rVert & \le\sum_{X\nsubseteq A}\lVert h_X\rVert\le2L.\notag
    \end{align}
\end{proof}

\begin{lemma}[Energy from ground-state overlap for a range-one chain]
    \label{thm:chain_range_one_overlap_energy}
    \lean{TNLean.MPS.AreaLaw.ChainHamiltonian.re_inner_sub_le_of_overlap_of_range_one,
        TNLean.MPS.AreaLaw.ChainHamiltonian.re_inner_le_of_overlap_div_of_range_one}
    \leanok
    \uses{def:chain_model}
    Let $H$ be a range-one Hamiltonian on a chain of $L$ sites, with each
    interaction term of norm at most one. Let $E_0,\Delta$ be real and let
    $\Omega$ be a unit vector satisfying $H\Omega=E_0\Omega$ and
    $H-E_0\ge\Delta(1-|\Omega\rangle\langle\Omega|)$.
    If $v$ is a unit vector with $|\langle\Omega,v\rangle|\ge1-\delta$, then
    \begin{align}
        \operatorname{Re}\langle v,Hv\rangle-E_0
         & \le8L\delta.
        \label{eq:chain_range_one_overlap_energy}
    \end{align}
    In particular, when $L>0$, overlap at least $1-c/L$ gives energy at
    most $E_0+8c$.

    This is a general norm estimate supporting the first energy inference
    in Lemma~l:within of~\cite[lines 509--515]{Landau2015Polynomial}.
    The source assumes positive adjacent-pair terms; here singleton and
    adjacent-pair terms need only be Hermitian. The correction for
    unsquared overlap is recorded in \cite{gap:lvv13_witness_energy}.
\end{lemma}
\begin{proof}\leanok
    \uses{thm:chain_range_one_support_count}
    The support count gives $\|H\|\le2L$. The eigenvector equation and
    $\|\Omega\|=1$ give $|E_0|\le\|H\|$, hence
    $\|H-E_0\,1\|\le\|H\|+|E_0|\le4L$.
    The energy estimate from overlap with an eigenvector in
    \cite[``Energy bound from overlap with an eigenvector'']{QICLean2026Blueprint}
    now gives (\ref{eq:chain_range_one_overlap_energy}). Substituting
    $\delta=c/L$ proves the last assertion.
\end{proof}

\begin{lemma}[Nearest-neighbour supports at two cuts]
    \label{thm:chain_two_cut_edge_supports}
    \lean{TNLean.MPS.AreaLaw.interval_edge_classification_of_lt}
    \leanok
    \uses{def:chain_model}
    Let $L,u,v$ be natural numbers with $u<v$ and $v+1<L$.
    Place the cuts immediately after sites $u$ and $v$, as in
    \cite{Arad2016RigorousRG}. Set $T=\{0,\ldots,u\}$, $E=\{u+1,\ldots,v\}$
    and $U=\{0,\ldots,L-1\}\setminus(T\cup E)$.
    For an edge $r$ with $r+1<L$, write $D_r=\{r,r+1\}$.
    Then $D_u\subseteq T\cup E$, $D_u\cap T=\{u\}$,
    $D_v\cap T=\varnothing$ and $D_v\cap E=\{v\}$.
    Every other edge has its support contained in $T$, in $E$, or in $U$.
    Adjacent cuts, $v=u+1$, are included.
\end{lemma}
\begin{proof}\leanok
    \uses{def:chain_model}
    Membership in $T\cup E$ is equivalent to having coordinate less than
    $v+1$; membership in $U$ is equivalent to having coordinate at least
    $v+1$. Comparing the endpoints gives the four assertions for $D_u,D_v$.
    For any other edge, either $r<u$, $u<r<v$, or $v<r$.
    These cases give $r+2\le u+1$, the pair of inequalities
    $u+1\le r$ and $r+2\le v+1$, or $v+1\le r$, respectively.
    The corresponding interval containment places $D_r$ in $T$, $E$, or $U$.
    For adjacent cuts, $E=\{u+1\}=\{v\}$. The first two assertions give
    $D_u\cap E=\{u+1\}$ and the fourth gives $D_v\cap E=\{v\}$, so the
    crossing supports $D_u$ and $D_v$ share the single middle site.
\end{proof}

\begin{lemma}[Changed sites under interval displacement]
    \label{thm:chain_interval_symmdiff}
    \lean{TNLean.MPS.AreaLaw.card_interval_symmDiff_le_dist}
    \leanok
    \uses{def:chain_model}
    For natural numbers $L,a,b,c,d$, let
    $[a,b)_L=\{t\in\{0,\ldots,L-1\}:a\le t<b\}$.
    Then
    \begin{align}
        |{\lbrack a,b)_L}\mathbin{\triangle}{\lbrack c,d)_L}| & \le |a-c|+|b-d|.
        \label{eq:chain_interval_symmdiff}
    \end{align}
    No ordering or containment assumptions on the endpoints are required.
\end{lemma}
\begin{proof}\leanok
    \uses{def:chain_model}
    A site belonging to exactly one interval lies between the two left
    endpoints or between the two right endpoints. Thus
    \begin{align}
        {\lbrack a,b)_L}\mathbin{\triangle}{\lbrack c,d)_L}
         & \subseteq [\min(a,c),\max(a,c))_L\notag \\
         & \qquad\cup [\min(b,d),\max(b,d))_L.
        \label{eq:chain_interval_endpoint_strips}
    \end{align}
    The two intervals in~(\ref{eq:chain_interval_endpoint_strips}) have at
    most $\max(a,c)-\min(a,c)=|a-c|$ and
    $\max(b,d)-\min(b,d)=|b-d|$ sites. Subadditivity of cardinality gives
    (\ref{eq:chain_interval_symmdiff}).
\end{proof}

\begin{lemma}[Middle-region displacement within two windows]
    \label{thm:chain_two_window_middle_displacement}
    \lean{TNLean.MPS.AreaLaw.card_two_window_middle_symmDiff_le}
    \leanok
    \uses{def:chain_model}
    Let $L,a,b,\ell$ be natural numbers, let $0\le p,s<2\ell$, and let
    $\delta\in\{0,1\}$. Using the clipped intervals of
    Lemma~\ref{thm:chain_interval_symmdiff}, set
    \begin{align}
        E        & =[a+p+1,b+s+1)_L,\notag      \\
        E_\delta & =[a+\ell+\delta,b+\ell+1)_L.
        \label{eq:chain_two_window_middle_regions}
    \end{align}
    Then $|E\mathbin{\triangle}E_\delta|\le2\ell$.
    When $a+\ell+1<b+\ell+1\le L$, write $i_1=a+\ell+1$ and
    $i_2=b+\ell+1$. In one-based site coordinates, $E_0$ is the block
    $\{i_1,\ldots,i_2\}$ in~\cite[Proposition~5.7]{Arad2016RigorousRG},
    whereas $E_1$ is the middle
    region between the cuts after $i_1$ and $i_2$ in its Theorem~4.2.
\end{lemma}
\begin{proof}\leanok
    \uses{thm:chain_interval_symmdiff}
    The inequalities $0\le p,s<2\ell$ and $\delta\in\{0,1\}$ give
    $|(p+1)-(\ell+\delta)|\le\ell$ and $|s-\ell|\le\ell$.
    Apply~(\ref{eq:chain_interval_symmdiff}) to the two intervals
    in~(\ref{eq:chain_two_window_middle_regions}) and add these bounds.
\end{proof}

\begin{lemma}[Frequency vectors with two prescribed coordinates]
    \label{thm:chain_two_fixed_frequencies}
    \lean{Nat.card_fun_sum_eq_and_two_eq,
        Nat.card_fun_sum_eq_and_two_eq_of_lt,
        Nat.card_fun_sum_eq_and_two_eq_of_total}
    \leanok
    Let $I$ be a finite set, let $u,v\in I$ be distinct, and let
    $r,i,j$ be natural numbers. Define
    \begin{align}
        \mathcal F_{r,i,j}
         & =\left\{\alpha:I\to\mathbb N:
        \sum_{a\in I}\alpha(a)=r,\quad
        \alpha(u)=i,\quad\alpha(v)=j\right\}.\label{eq:chain_fixed_frequency_set}
    \end{align}
    If $r<i+j$, then the set in~(\ref{eq:chain_fixed_frequency_set}) is empty.
    If $|I|=s+3$ for a natural number $s$, then
    \begin{align}
        |\mathcal F_{r,i,j}|
         & =\begin{cases}
                \displaystyle\binom{r-i-j+s}{s}, & i+j\le r, \\
                0,                               & r<i+j.
            \end{cases}\label{eq:chain_fixed_frequency_count}
    \end{align}
    Taking $r=m+i+j$ in~(\ref{eq:chain_fixed_frequency_count}) gives
    $|\mathcal F_{m+i+j,i,j}|=\binom{m+s}{s}$
    for every natural number $m$.
    Thus the count is one when $s=0$, provided $i+j\le r$,
    and also when $r=i+j$.
    The grouped alphabet in \cite[Section~4.1]{Arad2016RigorousRG}
    has two windows of $2\ell$ bond labels and three bulk labels;
    for this alphabet the formula applies with $s=4\ell$.
\end{lemma}
\begin{proof}\leanok
    For every $\alpha:I\to\mathbb N$,
    $\alpha(u)+\alpha(v)\le\sum_{a\in I}\alpha(a)$,
    proving the empty case.
    When the total is $m+i+j$, restricting to $I\setminus\{u,v\}$
    is a bijection between $\mathcal F_{m+i+j,i,j}$ and the
    nonnegative integer functions on $s+1$ labels whose sum is $m$.
    Its inverse reinserts $i$ at $u$ and $j$ at $v$.
    Identifying these functions with multisets of cardinality $m$
    gives $\binom{m+s}{s}$.
    When $i+j\le r$, substituting $m=r-i-j$ gives~(\ref{eq:chain_fixed_frequency_count}).
\end{proof}

\begin{lemma}[Uniform bound for frequency slices]
    \label{thm:chain_fixed_frequency_uniform_bound}
    \lean{Nat.card_fun_sum_eq_and_two_eq_le}
    \leanok
    Let $I$ have cardinality $s+3$, and let $u,v\in I$ be distinct.
    For natural numbers $r,i,j$, let $\mathcal F_{r,i,j}$ be the set
    in~(\ref{eq:chain_fixed_frequency_set}). Then
    \begin{align}
        |\mathcal F_{r,i,j}| & \le (r+1)^s.
        \label{eq:chain_fixed_frequency_uniform_bound}
    \end{align}
\end{lemma}
\begin{proof}\leanok
    \uses{thm:chain_two_fixed_frequencies}
    If $r<i+j$, the set is empty. Otherwise put $m=r-i-j$.
    The exact count in~(\ref{eq:chain_fixed_frequency_count}) and the
    elementary binomial estimate give
    \begin{align}
        |\mathcal F_{r,i,j}|=\binom{m+s}{s}
         & \le (m+1)^s\le(r+1)^s.
        \label{eq:chain_frequency_binomial_bound}
    \end{align}
    Combining~(\ref{eq:chain_frequency_binomial_bound}) with the empty case
    proves~(\ref{eq:chain_fixed_frequency_uniform_bound}).
\end{proof}

\begin{lemma}[Frequency slices through a prescribed degree]
    \label{thm:chain_fixed_frequency_degree_sum}
    \lean{Nat.sum_card_fun_sum_eq_and_two_eq_le}
    \leanok
    Let $I$ be a finite set of cardinality $s+3$, with $s\in\mathbb N$,
    and let $u,v\in I$ be distinct. For natural numbers $k,i,j$, use
    $\mathcal F_{r,i,j}$ from~(\ref{eq:chain_fixed_frequency_set}). Then
    \begin{align}
        \sum_{r=0}^{k}|\mathcal F_{r,i,j}|
         & \le (k+1)^{s+1}.
        \label{eq:chain_frequency_degree_sum}
    \end{align}
    All degrees from $0$ through $k$ are included. For an alphabet of
    $4\ell+3$ labels, the exponent is $4\ell+1$.
\end{lemma}
\begin{proof}\leanok
    \uses{thm:chain_fixed_frequency_uniform_bound}
    For each $0\le r\le k$,
    Lemma~\ref{thm:chain_fixed_frequency_uniform_bound} gives
    $|\mathcal F_{r,i,j}|\le(r+1)^s\le(k+1)^s$.
    Summing over the $k+1$ degrees gives
    $(k+1)(k+1)^s=(k+1)^{s+1}$.
\end{proof}

The following scalar estimates are auxiliaries for the two-window argument
of \cite[Proposition~5.7]{Arad2016RigorousRG}. For natural
$k,d,\ell,b$, write
\begin{align}
    C_{k,d,\ell,b}
     & =(2\ell)^2(b+1)^2(k+1)^{4\ell+1}\notag            \\
     & \qquad\cdot\bigl(2(k+1)\bigr)^{4b+2}d^{4b+4\ell}.
    \label{eq:chain_two_window_count_factor}
\end{align}
The full alphabet of $4\ell+3$ labels is retained, consisting of the two
windows and the three bulk labels.

\begin{lemma}[Summing the two-window counting factors]
    \label{thm:chain_two_window_count_sum}
    \lean{Nat.two_window_rank_sum_le_factor}
    \leanok
    Let $k,d,\ell,b\in\mathbb N$ with $d\ge1$ and $b\le k$. Then
    \begin{align}
         & (2\ell)^2
        \sum_{\substack{0\le r\le k  \\0\le i,j\le b}}
        (r+1)^{4\ell}
        \binom{r+i+1}{2i+1}\binom{r+j+1}{2j+1}
        d^{2(i+j)+4\ell}\notag       \\
         & \qquad\le C_{k,d,\ell,b}.
        \label{eq:chain_two_window_count_sum}
    \end{align}
    The sum ranges over every natural triple in the indicated ranges,
    including degree zero. The binomial factors retain the upper arguments
    $r+i+1$ and $r+j+1$ printed in
    \cite[Proposition~5.7]{Arad2016RigorousRG}; the count of commuting gaps
    is recorded in~\cite{gap:aklv13_interpolation_and_word_count}.
\end{lemma}
\begin{proof}\leanok
    For $r\le k$ and $i\le b\le k$, the elementary binomial estimate gives
    \begin{align}
        \binom{r+i+1}{2i+1}
         & \le (r+i+1)^{2i+1}
        \le\bigl(2(k+1)\bigr)^{2b+1}.
        \label{eq:chain_two_window_gap_bound}
    \end{align}
    Apply~(\ref{eq:chain_two_window_gap_bound}) to both binomials.
    Also $(r+1)^{4\ell}\le(k+1)^{4\ell}$ and
    $d^{2(i+j)+4\ell}\le d^{4b+4\ell}$, since $d\ge1$.
    Each summand is therefore at most
    $(k+1)^{4\ell}\bigl(2(k+1)\bigr)^{4b+2}d^{4b+4\ell}$.
    There are $(k+1)(b+1)^2$ triples. Multiplying by this number and by
    $(2\ell)^2$ gives the product
    (\ref{eq:chain_two_window_count_factor}).
\end{proof}

\begin{lemma}[An explicit bound for the counting product]
    \label{thm:chain_two_window_count_power}
    \lean{Nat.two_window_rank_factor_le_pow}
    \leanok
    Let $k,d,\ell,b\in\mathbb N$ with $k\ge1$, $d\ge2$,
    $\ell\ge1$ and $b\le k$. For the product in
    (\ref{eq:chain_two_window_count_factor}),
    \begin{align}
        C_{k,d,\ell,b}
         & \le(dk)^{17(\ell+b)}.
        \label{eq:chain_two_window_count_power}
    \end{align}
\end{lemma}
\begin{proof}\leanok
    Put $X=dk\ge2$. The hypotheses give $d\le X$,
    $b+1\le k+1\le X$, $2\ell\le2^\ell\le X^\ell$ and
    $2(k+1)\le X^2$. Substituting these bounds into
    (\ref{eq:chain_two_window_count_factor}) yields
    \begin{align}
        C_{k,d,\ell,b}
         & \le X^{10\ell+12b+7}
        \le X^{17(\ell+b)}.
        \label{eq:chain_two_window_count_exponent}
    \end{align}
    The last step uses $\ell\ge1$, which implies
    $10\ell+12b+7\le17(\ell+b)$.
\end{proof}

\begin{lemma}[Distance and spheres on a finite path]
    \label{lem:chain_path_metric_spheres}
    \lean{SimpleGraph.pathGraph_edist, SimpleGraph.pathGraph_card_sphere_le_two}
    \leanok
    Let $L$ be a natural number, and let $P_L$ be the graph on
    $\{0,\ldots,L-1\}$ joining consecutive coordinates. For vertices $x,y$ and
    every natural number $d$,
    \begin{align}
        d_{P_L}(x,y)                     & =|x-y|, \\
        \bigl|\{y:d_{P_L}(x,y)=d\}\bigr| & \le2.
    \end{align}
    Thus the sphere growth bound holds with coefficient $2$ and exponent $0$,
    independently of $L$. This concerns the contiguous chain setting of
    Section~1 of~\cite{TNLean2026OneDimensional}.
\end{lemma}
\begin{proof}
    \leanok
    Each edge changes the coordinate by one, so the displacement of a walk is
    at most its length. Taking the infimum over walks gives one inequality.
    In the other direction, successive coordinates between $x$ and $y$ stay
    inside the chain and give a walk of length $|x-y|$. A site at distance $d$
    from $x$ has coordinate $x-d$ or $x+d$; discarding coordinates outside the
    chain leaves at most two sites. At $d=0$ the candidates coincide. When
    $L=0$ there are no vertices, so the assertions are vacuous.
\end{proof}

\begin{lemma}[Regional forms of a nearest-neighbour bond array]
    \label{thm:chain_two_cut_edge_factors}
    \lean{TNLean.MPS.AreaLaw.exists_threeSplit_edge_factors}
    \leanok
    \uses{thm:chain_two_cut_edge_supports,lem:peps_frame_split_forms}
    Let $L,q,u,v$ be natural numbers with $u<v$ and $v+1<L$.
    For each bond index $0\le r<L-1$, let $H_r$ be a complex operator
    supported on $D_r=\{r,r+1\}$. Use the partition $T,E,U$ of
    Lemma~\ref{thm:chain_two_cut_edge_supports}, and write $\widetilde H_r$
    for $H_r$ in the three-region coordinates of
    Lemma~\ref{lem:peps_frame_split_forms}.
    There are arrays $H^{\mathrm L}_r,H^{\mathrm M}_r,H^{\mathrm R}_r$
    of operators on $\mathcal H_T,\mathcal H_E,\mathcal H_U$, respectively,
    and crossing factors $A_u$ on $\mathcal H_T\otimes\mathcal H_E$ and
    $B_v$ on $\mathcal H_E\otimes\mathcal H_U$ with the following properties.
    For every $r\ne u,v$,
    \begin{align}
        \widetilde H_r
         & =(H^{\mathrm L}_r\otimes\id_E)\otimes\id_U
        +(\id_T\otimes H^{\mathrm M}_r)\otimes\id_U\notag \\
         & \quad+\id_{TE}\otimes H^{\mathrm R}_r.
        \label{eq:chain_two_cut_ordinary_forms}
    \end{align}
    The two distinguished operators satisfy
    \begin{align}
        \widetilde H_u & =A_u\otimes\id_U,
                       & \operatorname{OSR}(A_u) & \le q^2,\notag \\
        \widetilde H_v & =\id_T\otimes B_v,
                       & \operatorname{OSR}(B_v) & \le q^2.
        \label{eq:chain_two_cut_crossing_forms}
    \end{align}
    All tensors are read in the same left-associated coordinates, using
    canonical reassociation for $\id_T\otimes B_v$.
    The arrays and crossing factors are chosen simultaneously for the given
    bond array and cuts. The local dimension may be zero, and the cuts may
    be adjacent. These are tensor forms for the local terms in
    \cite[Sections~4.1 and~5.2]{Arad2016RigorousRG}.
\end{lemma}
\begin{proof}\leanok
    \uses{thm:chain_two_cut_edge_supports,lem:peps_frame_split_forms,
        thm:peps_supported_left_factor_rank,thm:peps_supported_right_factor_rank}
    The support classification places every ordinary bond in $T$, $E$ or
    $U$. Apply the corresponding single-region part of
    Lemma~\ref{lem:peps_frame_split_forms}, putting zero in the other two
    array entries. Making these choices over the finite bond index set gives
    the arrays in~(\ref{eq:chain_two_cut_ordinary_forms}).

    For the first crossing, the classification gives
    $D_u\subseteq T\cup E$ and $D_u\cap T=\{u\}$.
    Theorem~\ref{thm:peps_supported_left_factor_rank} therefore supplies
    the same $A_u$ for its tensor identity and the bound
    $q^{2|D_u\cap T|}=q^2$.
    For the second crossing, $D_v\cap T=\varnothing$ and
    $D_v\cap E=\{v\}$, so
    Theorem~\ref{thm:peps_supported_right_factor_rank} supplies the same
    $B_v$ for its associated tensor identity and its bound
    $q^{2|D_v\cap E|}=q^2$.
    These are~(\ref{eq:chain_two_cut_crossing_forms}), including $q=0$.
\end{proof}

\begin{lemma}[Regional forms of two windows and three bulk operators]
    \label{thm:chain_two_window_factors}
    \lean{TNLean.MPS.AreaLaw.exists_threeSplit_two_window_factors}
    \leanok
    \uses{def:chain_model,lem:peps_frame_split_forms}
    Let $L,q,\ell,a,b$ be natural numbers with $a+2\ell\le b$ and
    $b+2\ell\le L-1$. Consider the $4\ell+3$ distinct labels
    \begin{align}
        \mathcal I
         & =\{x_r:0\le r<2\ell\}\sqcup\{y_r:0\le r<2\ell\}
        \sqcup\{\lambda,\mu,\rho\}.
        \label{eq:chain_two_window_alphabet}
    \end{align}
    For each $\alpha\in\mathcal I$, let $F_\alpha$ be a complex operator
    on the chain. Suppose that $F_{x_r}$ is supported on $[a+r,a+r+2)$
    and $F_{y_r}$ on $[b+r,b+r+2)$, and that the three independently
    given bulk operators satisfy
    \begin{align}
        F_\lambda & \text{ is supported on }[0,a+1),\notag       \\
        F_\mu     & \text{ is supported on }[a+2\ell,b+1),\notag \\
        F_\rho    & \text{ is supported on }[b+2\ell,L).
        \label{eq:chain_two_window_bulk_supports}
    \end{align}
    The labels remain distinct when operators coincide or vanish.
    Choose integers $0\le p,s<2\ell$, put $u=a+p$ and $v=b+s$, and set
    \begin{align}
        T & =[0,u+1),\notag   \\
        E & =[u+1,v+1),\notag \\
        U & =[v+1,L).
        \label{eq:chain_two_window_regions}
    \end{align}
    Write $\widetilde F_\alpha$ for $F_\alpha$ in the three-region
    coordinates of Lemma~\ref{lem:peps_frame_split_forms}.
    There are families $F^{\mathrm L}_\alpha,F^{\mathrm M}_\alpha,
        F^{\mathrm R}_\alpha$ of operators on $\mathcal H_T,\mathcal H_E,
        \mathcal H_U$, respectively, and operators $A$ on
    $\mathcal H_T\otimes\mathcal H_E$ and $B$ on
    $\mathcal H_E\otimes\mathcal H_U$, chosen simultaneously, such that
    for every $\alpha\ne x_p,y_s$,
    \begin{align}
        \widetilde F_\alpha
         & =(F^{\mathrm L}_\alpha\otimes\id_E)\otimes\id_U
        +(\id_T\otimes F^{\mathrm M}_\alpha)\otimes\id_U\notag \\
         & \quad+\id_{TE}\otimes F^{\mathrm R}_\alpha.
        \label{eq:chain_two_window_ordinary_forms}
    \end{align}
    The same $A$ and $B$ satisfy
    \begin{align}
        \widetilde F_{x_p}    & =A\otimes\id_U,\notag  \\
        \operatorname{OSR}(A) & \le q^2,\notag         \\
        \widetilde F_{y_s}    & =\id_T\otimes B,\notag \\
        \operatorname{OSR}(B) & \le q^2.
        \label{eq:chain_two_window_crossing_forms}
    \end{align}
    All tensors use the same left-associated coordinates, with canonical
    reassociation for $\id_T\otimes B$. The local dimension may be zero.
    When $a+2\ell=b$, the middle bulk support is the singleton $\{b\}$.
    This is a tensor decomposition for the two-window grouping in
    \cite[Sections~4.1 and~5.2]{Arad2016RigorousRG}.
\end{lemma}
\begin{proof}\leanok
    \uses{thm:chain_two_cut_edge_supports,lem:peps_frame_split_forms,
        thm:peps_supported_left_factor_rank,thm:peps_supported_right_factor_rank}
    The hypotheses give $u<v$ and $v+1<L$.
    Lemma~\ref{thm:chain_two_cut_edge_supports} places each window edge
    other than $u,v$ wholly in one of $T,E,U$.
    The three bulk supports satisfy
    \begin{align}
        {[0,a+1)}       & \subseteq T,\notag \\
        {[a+2\ell,b+1)} & \subseteq E,\notag \\
        {[b+2\ell,L)}   & \subseteq U.
        \label{eq:chain_two_window_bulk_containments}
    \end{align}
    Thus every operator other than $F_{x_p},F_{y_s}$ is supported in one
    region. Apply the corresponding single-region form of
    Lemma~\ref{lem:peps_frame_split_forms}, putting zero in the other two
    regional factors. Choosing these factors over $\mathcal I$ gives the
    three families in~(\ref{eq:chain_two_window_ordinary_forms}).

    The first crossing support is contained in $T\cup E$ and meets $T$
    in the single site $u$. Theorem~\ref{thm:peps_supported_left_factor_rank}
    supplies $A$ with both its tensor identity and the bound $q^2$.
    The second crossing support is disjoint from $T$ and meets $E$ in
    the single site $v$. Theorem~\ref{thm:peps_supported_right_factor_rank}
    supplies $B$ with both its tensor identity and the bound $q^2$.
    These are~(\ref{eq:chain_two_window_crossing_forms}), including $q=0$.
\end{proof}

\section{Locality on a contiguous chain}

\begin{lemma}[Lieb--Robinson estimates on a contiguous chain]
    \label{thm:chain_lieb_robinson}
    \lean{TNLean.MPS.AreaLaw.norm_heisenberg_commutator_le_chainDistance,
        TNLean.MPS.AreaLaw.norm_heisenberg_sub_siteExpectation_chainNeighborhood_le}
    \leanok
    \uses{def:chain_model,def:qc_site_expectation}
    Let $L,q,R$ be natural numbers. Let $K$ be a finite set of interaction
    labels, and for each $k\in K$ let $X_k$ be a nonempty support of diameter
    at most $R$ in the contiguous chain $I=\{0,\ldots,L-1\}$.
    Let $h_k$ be Hermitian operators supported on $X_k$, and suppose that
    $b_0\in\mathbb R$ satisfies
    \begin{align}
        \sum_{k:\,x\in X_k}\|h_k\|\le b_0
    \end{align}
    for every site $x$. Put $H=\sum_{k\in K}h_k$ and
    $\tau_t(A)=e^{itH}Ae^{-itH}$, and set
    \begin{align}
        C_0 & =2+4|b_0|e^{2R},\notag      \\
        v   & =2(R+1)|b_0|e^{2R}+1,\notag \\
        \mu & =\tfrac12,\qquad
        C'=\frac{2C_0e}{1-e^{-1/2}}.\notag
    \end{align}
    These constants are independent of $L,q$, the interaction family, and
    the observable support. For $S\subseteq I$ let
    $d(S,y)=\min_{x\in S}|x-y|$, with $d(\varnothing,y)=\infty$ and
    $e^{-\mu\infty}=0$. For every operator $A$ supported on $S$, every site
    $y$, every operator $B_y$ supported on $\{y\}$, and every real $t$,
    \begin{align}
        \|[\tau_t(A),B_y]\|
        \le C_0|S|\|A\|\|B_y\|e^{v|t|}e^{-d(S,y)/2}.
    \end{align}
    This estimate allows every natural local dimension $q$.
    If $q\ge1$, define, for every real $\ell\ge0$,
    $N_\ell(S)=\{y\in I:\exists x\in S,\ |x-y|\le\ell\}$.
    The normalized partial-trace expectation $E_{N_\ell(S)}$ then satisfies
    \begin{align}
        \|\tau_t(A)-E_{N_\ell(S)}(\tau_t(A))\|
        \le C'|S|^2\|A\|e^{v|t|-\ell/2}.
    \end{align}
    Both estimates allow $S=\varnothing$. This is the contiguous-chain case
    of the lemma ``Lieb--Robinson bound on a chain'' in Part~II
    of~\cite{TNLean2026OneDimensional}, as used for its Theorem~B;
    the distinction from the broader finite-domain reading is recorded
    in~\cite{gap:tnlean_one_dimensional_connected_domain_scope}.
\end{lemma}
\begin{proof}
    \leanok
    \uses{lem:chain_path_metric_spheres,thm:area_law_supported_identity,
        thm:qc_graph_lieb_robinson_finset,
        thm:qc_graph_neighborhood_localization}
    Use the path graph on $I$. Its distance is $|x-y|$, each $X_k$ has
    graph diameter at most $R$ and at most $R+1$ sites, and its nonemptiness
    permits a choice of an anchor in $X_k$. The matrix-entry characterization
    of support supplies the hypotheses on the supported algebras.
    Apply the arbitrary-support graph commutator estimate with decay
    parameter $1$ and support-size bound $R+1$. If $S$ is nonempty,
    its least distance to $y$ is a finite nonnegative integer $d$, and
    $e^{-d}\le e^{-d/2}$. For $S=\varnothing$, the factor $|S|$ makes
    the commutator bound zero without selecting a point of $S$.

    For localization assume $q\ge1$ and put $m=\lfloor\ell\rfloor$.
    Every distance sphere of the path has at most two sites, so the graph
    localization estimate applies with $K_g=2$ and growth exponent zero.
    Since path distances are integers, the union of graph balls of radius
    $m$ about $S$ is exactly $N_\ell(S)$. The tail constant and the rounding
    bound are
    \begin{align}
        T_0      & =\frac{e^{1/2}}{1-e^{-1/2}},\notag \\
        e^{-m/2} & \le e^{1/2}e^{-\ell/2}.\notag
    \end{align}
    Thus the graph coefficient $2C_0T_0$, multiplied by the rounding factor
    $e^{1/2}$, is $C'$. Empty unions and the zero factor $|S|^2$ cover
    $S=\varnothing$ as well.
\end{proof}
